%************************************************
\chapter{The Integration Challenge}
\label{ch:integration-challenge}
%************************************************

\section{Introduction}
\label{sec:integration-intro}

Systems biology seeks to understand living systems through \textbf{integration} of multiple biological scales:
\begin{itemize}
    \item Molecular interactions (protein-protein, enzyme-substrate)
    \item Metabolic networks (glycolysis, TCA cycle, respiration)
    \item Gene regulatory networks (transcription, translation, feedback)
    \item Signaling pathways (MAPK, PI3K/AKT, cAMP signaling)
    \item Cellular processes (cell cycle, apoptosis, differentiation)
\end{itemize}

\subsection{The Multi-Scale Challenge}
\label{subsec:multiscale-challenge}

These scales operate with \textbf{fundamentally different dynamics}:
\begin{itemize}
    \item \textbf{Fast}: Enzyme kinetics (milliseconds to seconds)
    \item \textbf{Medium}: Metabolic fluxes (seconds to minutes)
    \item \textbf{Slow}: Gene expression (minutes to hours)
    \item \textbf{Ultra-slow}: Cell cycle (hours to days)
\end{itemize}

\paragraph{Traditional approach:} Model each scale separately
\begin{itemize}
    \item Metabolic model: ODE system for enzyme kinetics
    \item Genetic model: Stochastic simulation for gene expression
    \item Signaling model: Boolean network or ODEs
\end{itemize}

\paragraph{Problem:} Real biology is \textbf{integrated}
\begin{itemize}
    \item Metabolite concentrations regulate gene expression (e.g., glucose represses lac operon)
    \item Proteins regulate metabolic fluxes (e.g., allosteric feedback inhibition)
    \item Gene expression produces enzymes that catalyze metabolism
    \item \textbf{Cross-scale regulation} is pervasive, not exceptional
\end{itemize}

\subsection{Chapter Overview}
\label{subsec:chapter3-overview}

This chapter:
\begin{itemize}
    \item Presents a \textbf{motivating example} (cAMP-CRP regulation of lac operon) demonstrating all four required capabilities
    \item Derives \textbf{eight formal requirements} (R1--R8) for integrated biological modeling
    \item Shows why \textbf{existing formalisms fail} to meet these requirements
    \item Motivates the \textbf{Extended Bio-Petri Net formalism} presented in Chapters 4--6
\end{itemize}

%************************************************
\section{Motivating Example: cAMP-CRP Regulation of Lac Operon}
\label{sec:lac-operon-example}
%************************************************

\subsection{Biological Background}
\label{subsec:lac-operon-background}

The \textbf{lac operon} in \textit{E. coli} is a classic model of gene regulation:

\paragraph{Components:}
\begin{itemize}
    \item \textbf{lacZ gene}: Encodes $\beta$-galactosidase (cleaves lactose $\to$ glucose + galactose)
    \item \textbf{lacY gene}: Encodes permease (imports lactose into cell)
    \item \textbf{lacI gene}: Encodes repressor (blocks transcription)
    \item \textbf{Promoter}: DNA binding site for RNA polymerase
    \item \textbf{Operator}: DNA binding site for repressor (overlaps promoter)
\end{itemize}

\paragraph{Regulation (simplified):}
\begin{enumerate}
    \item \textbf{No lactose}: Repressor binds operator $\to$ transcription blocked
    \item \textbf{Lactose present}: Lactose binds repressor $\to$ repressor releases $\to$ transcription ON
    \item \textbf{Glucose present}: cAMP low $\to$ CRP inactive $\to$ transcription LOW (even with lactose)
    \item \textbf{No glucose}: cAMP high $\to$ CRP-cAMP complex binds promoter $\to$ transcription HIGH
\end{enumerate}

\paragraph{Biological interpretation:}
\begin{itemize}
    \item \textbf{Lactose} is the \textbf{inducer} (removes repression)
    \item \textbf{Glucose} is the \textbf{preferred substrate} (via cAMP-CRP catabolite repression)
    \item \textbf{Logic}: Use glucose first (high ATP yield), then lactose if glucose depleted
    \item \textbf{Cross-scale regulation}:
    \begin{itemize}
        \item Metabolite (glucose) regulates signaling (cAMP)
        \item Signaling (cAMP-CRP) regulates transcription (lac genes)
        \item Transcription produces enzyme ($\beta$-galactosidase)
        \item Enzyme regulates metabolism (lactose $\to$ glucose)
    \end{itemize}
\end{itemize}

\subsection{Detailed Mechanistic Model}
\label{subsec:lac-model-detailed}

\paragraph{Species (11 places):}
\begin{enumerate}
    \item \textbf{Glucose} (metabolite, continuous concentration)
    \item \textbf{Lactose} (metabolite, continuous concentration)
    \item \textbf{cAMP} (signaling molecule, continuous)
    \item \textbf{CRP} (catabolite repressor protein, continuous)
    \item \textbf{CRP-cAMP} complex (active transcription factor, continuous)
    \item \textbf{lac\_gene} (DNA, discrete count = 1)
    \item \textbf{mRNA\_lac} (messenger RNA, discrete count, stochastic)
    \item \textbf{BetaGal} ($\beta$-galactosidase enzyme, continuous)
    \item \textbf{Repressor} (LacI protein, continuous)
    \item \textbf{Repressor-Lactose} complex (inactive repressor, continuous)
    \item \textbf{ATP} (energy currency, continuous)
\end{enumerate}

\paragraph{Transitions (9 reactions):}

\subsubsection{T1: Glucose Consumption (Continuous, Michaelis-Menten)}

\begin{equation}
\text{Glucose} + \text{ATP} \to \text{G6P} + \text{ADP}
\end{equation}

Rate:
\begin{equation}
v = \frac{V_{\max} \cdot [\text{Glucose}]}{K_m + [\text{Glucose}]}
\end{equation}

\begin{itemize}
    \item Consumes glucose (normal arc)
    \item Produces G6P (enters glycolysis)
    \item Continuous transition (ODE integration)
\end{itemize}

\subsubsection{T2: cAMP Synthesis (Continuous, Inhibited by Glucose)}

\begin{equation}
\text{ATP} \to \text{cAMP}
\end{equation}

Rate:
\begin{equation}
v = \frac{k_{\text{syn}}}{1 + [\text{Glucose}]/K_i}
\end{equation}

\begin{itemize}
    \item \textbf{Inhibitor arc}: Glucose $\multimap$ T2 (threshold: $K_i = 0.5$ mM)
    \item When $[\text{Glucose}] > K_i$: cAMP synthesis reduced
    \item When $[\text{Glucose}] < K_i$: cAMP synthesis increased
    \item \textbf{Cross-scale coupling}: Metabolite (glucose) regulates signaling (cAMP)
\end{itemize}

\subsubsection{T3: CRP-cAMP Complex Formation (Continuous, Mass Action)}

\begin{equation}
\text{CRP} + \text{cAMP} \rightleftharpoons \text{CRP-cAMP}
\end{equation}

Rates:
\begin{align}
v_f &= k_f \cdot [\text{CRP}] \cdot [\text{cAMP}] \quad \text{(forward)} \\
v_r &= k_r \cdot [\text{CRP-cAMP}] \quad \text{(reverse)}
\end{align}

\begin{itemize}
    \item Reversible reaction (two transitions, or single with negative flux)
    \item Fast equilibrium ($K_d = k_r / k_f \approx 1$ $\mu$M)
\end{itemize}

\subsubsection{T4: Repressor-Lactose Binding (Continuous, Fast Equilibrium)}

\begin{equation}
\text{Repressor} + \text{Lactose} \rightleftharpoons \text{Repressor-Lactose} \quad \text{(inactive repressor)}
\end{equation}

\begin{itemize}
    \item Lactose inactivates repressor (allosteric binding)
    \item Fast equilibrium
\end{itemize}

\subsubsection{T5: Lac Operon Transcription (Stochastic, Burst Mode)}

\begin{equation}
\text{lac\_gene} \to \text{lac\_gene} + \text{mRNA\_lac}
\end{equation}

\begin{itemize}
    \item \textbf{Test arc}: lac\_gene $\to$ T5 (gene not consumed, catalytic role)
    \item \textbf{Test arc}: CRP-cAMP $\to$ T5 (transcription factor enhances rate)
    \item \textbf{Inhibitor arc}: Repressor $\multimap$ T5 (repressor blocks transcription)
    \item \textbf{Stochastic transition}: Gillespie algorithm, propensity:
    \begin{equation}
    a = \begin{cases}
    k_{\text{basal}} + k_{\text{enhanced}} \cdot [\text{CRP-cAMP}] & \text{if } [\text{Repressor}] < \text{threshold} \\
    0 & \text{if } [\text{Repressor}] \geq \text{threshold}
    \end{cases}
    \end{equation}
    \item \textbf{Burst mode} ($\tau = \text{Burst}$): Produces 5--10 mRNA copies per event
    \item \textbf{Cross-scale coupling}: Signaling (cAMP-CRP) regulates genetics (transcription)
\end{itemize}

\subsubsection{T6: mRNA Translation (Stochastic, Discrete)}

\begin{equation}
\text{mRNA\_lac} \to \text{mRNA\_lac} + \text{BetaGal}
\end{equation}

\begin{itemize}
    \item \textbf{Test arc}: mRNA\_lac $\to$ T6 (mRNA template not consumed during translation)
    \item Stochastic transition: Each mRNA produces $\sim$100 proteins over lifetime
    \item Propensity: $a = k_{\text{translate}} \cdot [\text{mRNA\_lac}]$
\end{itemize}

\subsubsection{T7: mRNA Degradation (Stochastic, First-Order)}

\begin{equation}
\text{mRNA\_lac} \to \emptyset
\end{equation}

\begin{itemize}
    \item Exponential lifetime: $t_{\text{half}} \approx 3$ minutes
    \item Propensity: $a = k_{\text{deg}} \cdot [\text{mRNA\_lac}]$
\end{itemize}

\subsubsection{T8: Lactose Import (Continuous, Enzyme-Catalyzed)}

\begin{equation}
\text{Lactose}_{\text{ext}} \to \text{Lactose}
\end{equation}

\begin{itemize}
    \item \textbf{Test arc}: Permease $\to$ T8 (enzyme not consumed)
    \item For simplicity, assume Permease $\propto$ BetaGal (both from lac operon)
\end{itemize}

\subsubsection{T9: Lactose Cleavage (Continuous, Michaelis-Menten)}

\begin{equation}
\text{Lactose} \to \text{Glucose} + \text{Galactose}
\end{equation}

Rate:
\begin{equation}
v = \frac{V_{\max} \cdot [\text{BetaGal}] \cdot [\text{Lactose}]}{K_m + [\text{Lactose}]}
\end{equation}

\begin{itemize}
    \item \textbf{Test arc}: BetaGal $\to$ T9 (enzyme not consumed)
    \item \textbf{Cross-scale coupling}: Enzyme (from genetics) regulates metabolism
\end{itemize}

\subsection{System Dynamics}
\label{subsec:lac-dynamics}

\begin{table}[htbp]
\centering
\caption{Lac operon system dynamics over time}
\label{tab:lac-dynamics}
\begin{tabular}{lp{9cm}}
\toprule
\textbf{Phase} & \textbf{Description} \\
\midrule
\textbf{Initial state} ($t=0$) & 
$[\text{Glucose}] = 5$ mM (high, preferred substrate); 
$[\text{Lactose}] = 10$ mM (available, but unused);
$[\text{cAMP}] = 0.1$ $\mu$M (low, due to glucose);
$[\text{CRP-cAMP}] \approx 0$;
$[\text{mRNA\_lac}] = 0$ (transcription repressed);
$[\text{BetaGal}] = 10$ molecules (basal level) \\
\midrule
\textbf{Phase 1} ($t=0$ to $20$ min) & 
T1 fires continuously: Glucose $\to$ G6P (glycolysis);
$[\text{Glucose}]$ decreases: $5 \to 0.1$ mM;
T2 (cAMP synthesis) gradually activated as glucose drops;
$[\text{cAMP}]$ increases: $0.1 \to 10$ $\mu$M (100-fold);
T3: CRP + cAMP $\to$ CRP-cAMP;
$[\text{CRP-cAMP}]$ increases: $0 \to 5$ $\mu$M \\
\midrule
\textbf{Phase 2} ($t=20$ to $40$ min) & 
$[\text{CRP-cAMP}]$ high $\to$ T5 (transcription) activated;
T5 fires stochastically: Bursts of 5--10 mRNA copies;
$[\text{mRNA\_lac}]$ increases: $0 \to 50$ molecules (stochastic fluctuations);
T6 fires stochastically: mRNA $\to$ BetaGal;
$[\text{BetaGal}]$ increases: $10 \to 5000$ molecules (500-fold induction) \\
\midrule
\textbf{Phase 3} ($t=40$ to $80$ min) & 
$[\text{BetaGal}]$ high $\to$ T9 activated;
T9 fires continuously: Lactose $\to$ Glucose + Galactose;
$[\text{Lactose}]$ decreases: $10 \to 0$ mM;
$[\text{Glucose}]$ increases: $0.1 \to 5$ mM (replenished from lactose);
\textbf{Feedback}: Glucose rise $\to$ $[\text{cAMP}]$ drops $\to$ transcription reduced;
System returns to low basal state \\
\midrule
\textbf{Steady state} ($t > 80$ min) & 
$[\text{Glucose}] = 5$ mM (maintained by lactose cleavage);
$[\text{Lactose}] = 0$ mM (depleted);
$[\text{BetaGal}] = 5000$ molecules (stable, long protein half-life);
$[\text{mRNA\_lac}] = 5$ molecules (basal level, transcription off) \\
\bottomrule
\end{tabular}
\end{table}

\subsection{Four Capabilities Required}
\label{subsec:lac-capabilities}

\subsubsection{Capability 1: Weak Independence \& Cooperative Parallelism}

\textbf{Example:} Multiple pathways producing ATP
\begin{itemize}
    \item Glycolysis: Glucose $\to$ ATP (via $t_1'$)
    \item Oxidative phosphorylation: \ce{O2} $\to$ ATP (via $t_{\text{oxphos}}$)
\end{itemize}

Properties:
\begin{itemize}
    \item Both produce ATP (shared output)
    \item Disjoint inputs: Glucose $\neq$ \ce{O2}
    \item \textbf{Weakly independent}: Can compute both fluxes in parallel, effects add to ATP pool
\end{itemize}

\textbf{In lac operon:} T5 (transcription) and T6 (translation) both interact with mRNA\_lac
\begin{itemize}
    \item T5 produces mRNA (output)
    \item T6 uses mRNA as template (test arc, non-consumptive)
    \item \textbf{Weak independence}: YES (T6 does not consume mRNA)
    \item Can execute in parallel (superposition: multiple ribosomes translate same mRNA)
\end{itemize}

\subsubsection{Capability 2: Heterogeneous Transition Types Coexistence}

The lac operon model requires \textbf{four transition types} in a single network:

\begin{table}[htbp]
\centering
\caption{Heterogeneous transition types in lac operon}
\label{tab:transition-types}
\begin{tabular}{llp{6cm}}
\toprule
\textbf{Type} & \textbf{Transitions} & \textbf{Characteristics} \\
\midrule
Continuous & T1, T2, T3, T9 & ODE integration: $dM/dt = \Phi(M)$; Fast equilibrium reactions (T3) use stiff ODE solvers \\
Stochastic & T5, T6, T7 & Gillespie algorithm: Discrete events, exponentially distributed intervals; Low copy numbers (mRNA: 0--50 molecules) \\
Burst & T5 & Transcription occurs in pulses (RNA polymerase processivity); Each burst produces 5--10 mRNA copies (geometric distribution) \\
\bottomrule
\end{tabular}
\end{table}

\paragraph{Challenge:} Synchronize all types
\begin{itemize}
    \item Continuous ODE step: $\Delta t = 0.01$ s (adaptive)
    \item Stochastic Gillespie: $\tau_{\text{next}} = -\ln(\text{rand}()) / a_{\text{total}}$ (event-driven)
    \item Burst events: Independent exponential process
    \item \textbf{Hybrid scheduler} required (see Chapter 11)
\end{itemize}

\subsubsection{Capability 3: Arc-Level Regulation with Biochemical Semantics}

Regulatory logic is \textbf{embedded in the topology} (not hidden in code):

\paragraph{Test arcs (non-consumptive catalysis):}
\begin{itemize}
    \item lac\_gene $\to$ T5: Gene is template for transcription (not consumed)
    \item CRP-cAMP $\to$ T5: Transcription factor enhances rate (not consumed)
    \item mRNA\_lac $\to$ T6: mRNA is template for translation (not consumed)
    \item BetaGal $\to$ T9: Enzyme catalyzes reaction (not consumed)
\end{itemize}

\paragraph{Inhibitor arcs (threshold-based blocking):}
\begin{itemize}
    \item Glucose $\multimap$ T2: High glucose inhibits cAMP synthesis
    \begin{itemize}
        \item Threshold formula: $\Delta(\text{Glucose}, \text{T2}) = K_i = 0.5$ mM
        \item Enabling condition: $[\text{Glucose}] < 0.5$ mM
    \end{itemize}
    \item Repressor $\multimap$ T5: Repressor blocks transcription
    \begin{itemize}
        \item Threshold: $\Delta(\text{Repressor}, \text{T5}) = 10$ nM
        \item Enabling condition: $[\text{Repressor}] < 10$ nM
    \end{itemize}
\end{itemize}

\begin{figure}[htbp]
\centering
\includegraphics[width=0.8\textwidth]{gfx/lac_operon.pdf}
\caption{Lac operon network topology showing arc-level regulation. Solid arcs: normal (consumptive); Dashed green arcs: test (catalytic); Dotted red arcs: inhibitor (regulatory).}
\label{fig:lac-topology}
\end{figure}

% Original TikZ code archived for reference:
% \begin{tikzpicture}[...]
% (TikZ code removed - figure now generated from SHYpn model)
% \end{tikzpicture}

\paragraph{Key insight:} The entire regulatory logic is visible in the \textbf{network topology}. An external viewer can understand:
\begin{itemize}
    \item Glucose inhibits cAMP synthesis (inhibitor arc)
    \item CRP-cAMP activates transcription (test arc)
    \item Repressor blocks transcription (inhibitor arc)
    \item BetaGal catalyzes lactose cleavage (test arc)
    \item Glucose production feeds back to cAMP inhibition (cycle)
\end{itemize}

\subsubsection{Capability 4: Atomic Conservation \& Biochemical Formula Tracking}

Each place has \textbf{elemental composition}:
\begin{itemize}
    \item \textbf{Glucose}: \ce{C6H12O6}
    \item \textbf{Lactose}: \ce{C12H22O11} (disaccharide: Glucose + Galactose -- \ce{H2O})
    \item \textbf{cAMP}: \ce{C10H12N5O6P}
    \item \textbf{ATP}: \ce{C10H16N5O13P3}
\end{itemize}

\paragraph{Elemental balance verification:}

\textbf{T9: Lactose cleavage}
\begin{equation}
\ce{C12H22O11 + H2O -> C6H12O6 + C6H12O6}
\end{equation}

Balance check:
\begin{align}
\text{\ce{C}}: & \quad 12 = 6 + 6 \quad \checkmark \\
\text{\ce{H}}: & \quad 22 + 2 = 12 + 12 \quad \checkmark \quad (24 = 24) \\
\text{\ce{O}}: & \quad 11 + 1 = 6 + 6 \quad \checkmark \quad (12 = 12)
\end{align}

\textbf{T1: Glucose phosphorylation}
\begin{equation}
\ce{C6H12O6 + C10H16N5O13P3 -> C6H11O9P + C10H16N5O10P2 + H+}
\end{equation}

Balance check:
\begin{align}
\text{\ce{C}}: & \quad 6 + 10 = 6 + 10 \quad \checkmark \quad (16 = 16) \\
\text{\ce{H}}: & \quad 12 + 16 = 11 + 16 + 1 \quad \checkmark \quad (28 = 28) \\
\text{\ce{O}}: & \quad 6 + 13 = 9 + 10 \quad \checkmark \quad (19 = 19) \\
\text{\ce{N}}: & \quad 0 + 5 = 0 + 5 \quad \checkmark \quad (5 = 5) \\
\text{\ce{P}}: & \quad 0 + 3 = 1 + 2 \quad \checkmark \quad (3 = 3)
\end{align}

\paragraph{Purpose:} Elemental balance ensures \textbf{mass conservation} at the atomic level, not just token counting.

%************************************************
\section{Eight Formal Requirements for Integrated Biological Modeling}
\label{sec:eight-requirements}
%************************************************

Based on the lac operon example and broader analysis of systems biology models, we derive \textbf{eight formal requirements} (R1--R8):

\subsection{R1: Multi-Scale Temporal Dynamics (Milliseconds to Hours)}
\label{subsec:r1-multiscale}

\paragraph{Requirement:} Formalism must support processes spanning 6+ orders of magnitude in timescale.

\paragraph{Justification:}
\begin{itemize}
    \item Enzyme kinetics: $10^{-3}$ -- $10^0$ seconds (milliseconds to seconds)
    \item Metabolic steady state: $10^0$ -- $10^2$ seconds (seconds to minutes)
    \item Gene expression: $10^2$ -- $10^4$ seconds (minutes to hours)
    \item Cell cycle: $10^4$ -- $10^5$ seconds (hours to days)
\end{itemize}

\paragraph{Example (lac operon):}
\begin{itemize}
    \item cAMP-CRP binding: $t_{\text{equilibrium}} \approx 0.1$ s (fast)
    \item mRNA half-life: $t_{\text{half}} \approx 180$ s (3 minutes)
    \item Lac operon induction: $t_{\text{response}} \approx 1200$ s (20 minutes)
\end{itemize}

\paragraph{Implication:} Single timescale models (e.g., fixed $\Delta t$ ODE) cannot capture both fast equilibria and slow gene expression.

\paragraph{Solution:} Heterogeneous transition types (Continuous for fast, Stochastic for slow, adaptive time steps).

\subsection{R2: Discrete and Continuous State Variables (Hybrid Semantics)}
\label{subsec:r2-hybrid}

\paragraph{Requirement:} Formalism must support both discrete counts (genes, mRNA copies) and continuous concentrations (metabolites, proteins).

\paragraph{Justification:}
\begin{itemize}
    \item Genes: Discrete (1 copy per cell for bacterial chromosome, 2 copies for diploid)
    \item mRNA: Discrete (0--100 copies, stochastic fluctuations)
    \item Proteins: Continuous (1000--1,000,000 copies, law of large numbers applies)
    \item Metabolites: Continuous (millimolar concentrations, millions of molecules)
\end{itemize}

\paragraph{Example (lac operon):}
\begin{itemize}
    \item lac\_gene: Discrete (1 copy)
    \item mRNA\_lac: Discrete (0--50 copies, Poisson-like distribution)
    \item BetaGal: Continuous (10--10,000 molecules, treat as concentration)
    \item Glucose: Continuous (0--10 mM = $10^{16}$ molecules per cell, deterministic)
\end{itemize}

\paragraph{Implication:} Pure discrete models (e.g., Gillespie for entire network) are computationally intractable for high-copy metabolites. Pure continuous models (ODEs) miss stochastic gene expression noise.

\paragraph{Solution:} Hybrid semantics (discrete marking for low-copy, continuous marking for high-copy).

\subsection{R3: Non-Consumptive Participation (Catalysis, Templating)}
\label{subsec:r3-nonconsumptive}

\paragraph{Requirement:} Formalism must distinguish between consumptive reactions (substrates) and non-consumptive roles (enzymes, genes, transcription factors).

\paragraph{Justification:}
\begin{itemize}
    \item Enzymes: Participate in reactions but are not consumed ($k_{\text{cat}}$ turnover, enzyme released)
    \item Genes: Serve as templates for transcription (DNA not depleted)
    \item mRNA: Templates for translation (mRNA produces many proteins before degradation)
    \item Transcription factors: Bind DNA, enhance transcription, but are not degraded
\end{itemize}

\paragraph{Example (lac operon):}
\begin{itemize}
    \item lac\_gene $\to$ Transcription: Gene is template (test arc, non-consumptive)
    \item BetaGal $\to$ Lactose\_cleavage: Enzyme catalyzes (test arc, non-consumptive)
    \item If modeled with normal arcs: Gene and enzyme would deplete after single firing \texttimes
\end{itemize}

\paragraph{Implication:} Classical Petri nets (only normal arcs) cannot model catalysis without artificial workarounds (e.g., enzyme $\to$ enzyme + product, doubling token count).

\paragraph{Solution:} Test arcs ($p \to t$) with semantics: $M'(p) = M(p)$ (unchanged after firing).

\subsection{R4: Threshold-Based Regulation (Inhibition, Activation Thresholds)}
\label{subsec:r4-threshold}

\paragraph{Requirement:} Formalism must support enabling conditions based on \textbf{place markings exceeding or falling below thresholds}.

\paragraph{Justification:}
\begin{itemize}
    \item Allosteric inhibition: Product inhibits enzyme when $[\text{Product}] > K_i$ (e.g., ATP inhibits PFK)
    \item Transcriptional repression: Repressor blocks transcription when $[\text{Repressor}] > \text{threshold}$
    \item Ultrasensitivity: Sigmoidal response curves (Hill equation) require threshold-like behavior
\end{itemize}

\paragraph{Example (lac operon):}
\begin{itemize}
    \item Glucose $\multimap$ cAMP\_synthesis: Transition blocked when $[\text{Glucose}] > 0.5$ mM
    \item Repressor $\multimap$ Transcription: Transition blocked when $[\text{Repressor}] > 10$ nM
    \item CRP-cAMP $\to$ Transcription: Rate enhanced when $[\text{CRP-cAMP}] > \text{threshold}$ (could use Hill equation)
\end{itemize}

\paragraph{Implication:} Classical Petri nets have binary enabling ($M(p) \geq W(p,t)$ or not). Cannot express ``fire only if $M(p) < \text{threshold}$''.

\paragraph{Solution:} Inhibitor arcs ($p \multimap t$) with threshold functions $\Delta(p,t)$. Enabling condition: $M(p) < \Delta(p,t)$.

\subsection{R5: Cooperative and Convergent Processes (Shared Outputs, Superposition)}
\label{subsec:r5-cooperative}

\paragraph{Requirement:} Formalism must allow multiple transitions to \textbf{produce to the same place} without conflict (additive effects).

\paragraph{Justification:}
\begin{itemize}
    \item Convergent metabolism: Multiple pathways produce same metabolite (e.g., pyruvate from glycolysis, lactate, alanine)
    \item Parallel enzyme activity: Multiple enzyme copies act on same substrate pool (effects add)
    \item Multiple transcription factors: Several activators enhance same promoter (additive or synergistic)
\end{itemize}

\paragraph{General example (pyruvate):}
\begin{align}
\text{Glycolysis} &\to \text{Pyruvate} \quad \text{(flux } v_1\text{)} \\
\text{Lactate\_oxidation} &\to \text{Pyruvate} \quad \text{(flux } v_2\text{)} \\
\text{Alanine\_transamination} &\to \text{Pyruvate} \quad \text{(flux } v_3\text{)} \\
\text{Net: } \frac{dM(\text{Pyruvate})}{dt} &= v_1 + v_2 + v_3 \quad \text{(superposition)}
\end{align}

\paragraph{Implication:} Classical strong independence rejects transitions with shared outputs (conflict condition). But biologically, shared outputs are \textbf{cooperative}, not conflicting.

\paragraph{Solution:} Weak independence (disjoint inputs, shared outputs allowed). Parallel execution with synchronized accumulation.

\subsection{R6: Multi-Type Kinetics}
\label{subsec:r6-kinetics}

\paragraph{Requirement:} Formalism must support diverse rate functions within a single model:
\begin{itemize}
    \item Mass action, Michaelis-Menten, Hill, stochastic propensities
\end{itemize}

\paragraph{Justification:}
\begin{itemize}
    \item Fast binding: Mass action ($k \cdot [A] \cdot [B]$)
    \item Enzyme kinetics: Michaelis-Menten ($V_{\max} \cdot [S] / (K_m + [S])$)
    \item Cooperative regulation: Hill equation ($V_{\max} \cdot [S]^n / (K^n + [S]^n)$)
    \item Stochastic gene expression: Propensity functions ($a = k_{\text{basal}} + k_{\text{induced}} \cdot [\text{TF}]$)
    \item Bursts: Geometric distribution for burst size, exponential inter-burst intervals
\end{itemize}

\paragraph{Example (lac operon):}
\begin{table}[htbp]
\centering
\caption{Kinetic types in lac operon model}
\label{tab:lac-kinetics}
\begin{tabular}{llp{5cm}}
\toprule
\textbf{Transition} & \textbf{Kinetic Type} & \textbf{Rate Function} \\
\midrule
T1 & Michaelis-Menten & $V_{\max} \cdot [\text{Glc}] / (K_m + [\text{Glc}])$ \\
T2 & Inhibited synthesis & $k / (1 + [\text{Glc}]/K_i)$ \\
T3 & Mass action & $k_f \cdot [\text{CRP}] \cdot [\text{cAMP}]$ \\
T5 & Stochastic + burst & Propensity + geometric burst size \\
T9 & Michaelis-Menten & $V_{\max} \cdot [E] \cdot [S] / (K_m + [S])$ \\
\bottomrule
\end{tabular}
\end{table}

\paragraph{Implication:} Formalism must provide \textbf{rate function library} $\Phi$ and allow \textbf{per-transition customization}.

\paragraph{Solution:} Rate function map $\Phi: T \to \text{RateFunction}$. Each transition has type $\tau(t)$ and associated kinetic parameters.

\subsection{R7: Elemental Conservation (Mass Balance at Atomic Level)}
\label{subsec:r7-conservation}

\paragraph{Requirement:} Formalism must track \textbf{elemental composition} (\ce{C}, \ce{H}, \ce{O}, \ce{N}, \ce{P}, \ce{S}) and verify conservation laws.

\paragraph{Justification:}
\begin{itemize}
    \item Fundamental physics: Atoms are neither created nor destroyed (non-nuclear biochemistry)
    \item Debugging: Elemental imbalance indicates modeling error (missing cofactor, wrong stoichiometry)
    \item Redox balance: Electron conservation in oxidation-reduction reactions
    \item ATP budget: Phosphate balance in energy metabolism
\end{itemize}

\paragraph{Example (lac operon):}
\begin{itemize}
    \item Lactose (\ce{C12H22O11}) $\to$ Glucose (\ce{C6H12O6}) + Galactose (\ce{C6H12O6})
    \item Check: $12 = 6+6$ (\ce{C}), $22+2 = 12+12$ (\ce{H}), $11+1 = 6+6$ (\ce{O}) \checkmark
    \item If forgot water: $12=12$ (\ce{C}) \checkmark, but $22 \neq 24$ (\ce{H}) \texttimes\ $\to$ Error detected
\end{itemize}

\paragraph{Implication:} Classical Petri nets use abstract tokens (no chemical identity). Cannot verify conservation.

\paragraph{Solution:} Biochemical formula map $\rho: T \to \text{Formula}$. Elemental balance matrix $S_e$. Verify: $S_e \cdot v = 0$ (steady state).

\subsection{R8: Parallelism for Computational Efficiency}
\label{subsec:r8-parallelism}

\paragraph{Requirement:} Formalism must enable \textbf{parallel execution} of independent processes to achieve scalable performance.

\paragraph{Justification:}
\begin{itemize}
    \item Large models: Genome-scale metabolism (1000+ reactions), whole-cell models (5000+ processes)
    \item Real-time simulation: Parameter sweeps (100+ conditions), sensitivity analysis (1000+ samples)
    \item Multi-core hardware: Modern CPUs (8--64 cores), GPUs (1000+ cores)
\end{itemize}

\paragraph{Example (genome-scale):}
\begin{itemize}
    \item 1000 transitions, 650 weakly independent pairs
    \item Classical strong independence: 200 pairs $\to$ 5 parallel sets (speedup: 5$\times$)
    \item Weak independence: 650 pairs $\to$ 3 parallel sets (speedup: 10$\times$)
    \item \textbf{Result}: 2$\times$ speedup improvement by recognizing biological cooperativity
\end{itemize}

\paragraph{Implication:} Formalism must provide \textbf{dependency analysis} (which transitions conflict) and \textbf{parallel execution semantics} (safe concurrent firing).

\paragraph{Solution:} Weak independence theory (Chapter 5). Dependency classification algorithm. Parallel scheduler.

%************************************************
\section{Why Existing Formalisms Fail}
\label{sec:existing-formalisms}
%************************************************

\subsection{Comparative Analysis}
\label{subsec:comparative-analysis}

\begin{table}[htbp]
\centering
\caption{Existing formalisms vs.\ eight requirements}
\label{tab:formalism-comparison}
\small
\begin{tabular}{lcccccccc}
\toprule
\textbf{Formalism} & \textbf{R1} & \textbf{R2} & \textbf{R3} & \textbf{R4} & \textbf{R5} & \textbf{R6} & \textbf{R7} & \textbf{R8} \\
\midrule
Classical PN & \texttimes & \texttimes & \texttimes & \texttimes & $\sim$ & \texttimes & \texttimes & $\sim$ \\
Stochastic PN & \texttimes & $\sim$ & \texttimes & \texttimes & $\sim$ & $\sim$ & \texttimes & \texttimes \\
Continuous PN & \texttimes & $\sim$ & \texttimes & \texttimes & \checkmark & \checkmark & \texttimes & $\sim$ \\
Hybrid PN & $\sim$ & \checkmark & \texttimes & \texttimes & \checkmark & $\sim$ & \texttimes & $\sim$ \\
Colored PN & \texttimes & $\sim$ & \texttimes & $\sim$ & $\sim$ & \texttimes & $\sim$ & $\sim$ \\
Rule-Based & $\sim$ & \texttimes & $\sim$ & $\sim$ & \texttimes & $\sim$ & \texttimes & \texttimes \\
ODE Systems & \texttimes & \texttimes & $\sim$ & $\sim$ & \checkmark & \checkmark & $\sim$ & $\sim$ \\
\midrule
\textbf{Extended Bio-PN} & \checkmark & \checkmark & \checkmark & \checkmark & \checkmark & \checkmark & \checkmark & \checkmark \\
\bottomrule
\end{tabular}
\end{table}

\noindent\textbf{Legend:} \checkmark\ Fully addressed; $\sim$ Partially addressed; \texttimes\ Not addressed

\subsection{Classical Petri Nets (Place-Transition Nets)}
\label{subsec:classical-pn}

\paragraph{Definition:} 5-tuple $(P, T, F, W, M_0)$
\begin{itemize}
    \item Only normal arcs (consumptive)
    \item Homogeneous transitions (no types)
    \item Abstract tokens (no chemical identity)
\end{itemize}

\paragraph{Verdict:} Classical Petri nets provide the \textbf{topological foundation} (bipartite graph, firing rules) but lack \textbf{biochemical semantics} and \textbf{hybrid dynamics}.

\subsection{Continuous Petri Nets (CPNs)}
\label{subsec:continuous-pn}

\paragraph{Definition:} Classical PN + continuous marking ($M: P \to \mathbb{R}_{\geq 0}$) + rate functions $\Phi$
\begin{itemize}
    \item Transitions fire continuously (ODE: $dM/dt = C \cdot \Phi(M)$)
    \item Incidence matrix $C$
\end{itemize}

\paragraph{Verdict:} CPNs are excellent for \textbf{metabolic networks} (continuous, deterministic) but cannot model \textbf{stochastic gene expression} or \textbf{discrete molecular counts}.

\subsection{Hybrid Petri Nets (HPNs)}
\label{subsec:hybrid-pn}

\paragraph{Definition:} Classical PN + discrete places ($P_d$) + continuous places ($P_c$) + hybrid transitions
\begin{itemize}
    \item Combines discrete and continuous semantics
    \item Hybrid firing rules (discrete $\to$ discrete, continuous $\to$ continuous, discrete $\to$ continuous)
\end{itemize}

\paragraph{Verdict:} HPNs address \textbf{hybrid state} but lack \textbf{stochastic transitions}, \textbf{test/inhibitor arcs}, and \textbf{biochemical semantics}.

\subsection{ODE Systems (SBML Models)}
\label{subsec:ode-systems}

\paragraph{Definition:} Systems of ordinary differential equations:
\begin{equation}
\frac{dM(\text{Glucose})}{dt} = -v_{\text{HK}} + v_{\text{import}}, \quad \frac{dM(\text{G6P})}{dt} = v_{\text{HK}} - v_{\text{PGI}}, \quad \ldots
\end{equation}

\paragraph{Verdict:} ODE systems are the \textbf{standard} for metabolic modeling but lack \textbf{stochastic gene expression}, \textbf{topological visibility} of regulation, and \textbf{automatic elemental balance}.

%************************************************
\section{The Integration Gap}
\label{sec:integration-gap}
%************************************************

\subsection{Current Practice: Separate Models}
\label{subsec:separate-models}

\textbf{Standard approach} in systems biology:
\begin{enumerate}
    \item \textbf{Metabolic model}: SBML ODE system (continuous, deterministic)
    \item \textbf{Gene regulatory model}: Boolean network or stochastic simulation (discrete, probabilistic)
    \item \textbf{Signaling model}: ODE system or rule-based (Kappa)
\end{enumerate}

\paragraph{Integration attempt:}
\begin{itemize}
    \item Run metabolic model $\to$ compute $[\text{Glucose}]$
    \item Use $[\text{Glucose}]$ as input to gene regulatory model $\to$ compute $[\text{mRNA\_lac}]$
    \item Use $[\text{mRNA\_lac}]$ to update enzyme levels in metabolic model $\to$ update $[\text{BetaGal}]$
    \item \textbf{Problem}: Loose coupling, manual synchronization, no formal semantics
\end{itemize}

\paragraph{Problems:}
\begin{enumerate}
    \item \textbf{Manual synchronization}: Requires custom glue code for every model pair
    \item \textbf{No formal semantics}: When to synchronize? What if Gillespie event occurs mid-ODE step?
    \item \textbf{Error-prone}: Easy to miss feedback loops (e.g., Glucose $\to$ cAMP $\to$ Transcription $\to$ BetaGal $\to$ Lactose $\to$ Glucose)
    \item \textbf{Not reusable}: Custom integration script per model
    \item \textbf{No topological visibility}: Regulation hidden in code, not visible in network
\end{enumerate}

\subsection{The Need for Unified Formalism}
\label{subsec:unified-formalism}

\paragraph{Goal:} Single model representation capturing all scales and semantics.

\paragraph{Requirements for unified formalism:}
\begin{enumerate}
    \item \textbf{Single network}: All species (genes, mRNA, proteins, metabolites) as nodes
    \item \textbf{Hybrid dynamics}: Discrete (genes, mRNA) + continuous (metabolites)
    \item \textbf{Heterogeneous transitions}: Continuous (metabolism) + stochastic (gene expression) + burst (transcription)
    \item \textbf{Topological regulation}: Inhibitor/test arcs visible in network
    \item \textbf{Automatic synchronization}: Built-in hybrid scheduler
    \item \textbf{Elemental conservation}: Automatic verification
\end{enumerate}

\paragraph{Extended Bio-PN achieves this} (see Chapters 4--6).

%************************************************
\section{Summary}
\label{sec:chapter3-summary}
%************************************************

This chapter established the \textbf{integration challenge} in systems biology:

\paragraph{Key Contributions:}
\begin{enumerate}
    \item \textbf{Motivating example}: cAMP-CRP regulation of lac operon demonstrates:
    \begin{itemize}
        \item Cross-scale regulation (metabolite $\to$ signaling $\to$ genetics $\to$ metabolism)
        \item Heterogeneous dynamics (continuous metabolism + stochastic transcription + burst mode)
        \item Diverse arc types (test arcs for catalysis, inhibitor arcs for regulation)
        \item Elemental conservation (glucose, lactose, ATP formulas)
    \end{itemize}
    
    \item \textbf{Eight formal requirements} (R1--R8):
    \begin{itemize}
        \item R1: Multi-scale temporal dynamics
        \item R2: Hybrid discrete-continuous state
        \item R3: Non-consumptive participation (catalysis)
        \item R4: Threshold-based regulation (inhibition)
        \item R5: Cooperative and convergent processes (shared outputs)
        \item R6: Multi-type kinetics (Michaelis-Menten, Hill, stochastic)
        \item R7: Elemental conservation (mass balance)
        \item R8: Parallelism for computational efficiency
    \end{itemize}
    
    \item \textbf{Existing formalisms fail}:
    \begin{itemize}
        \item Classical Petri nets: Topological foundation, but lack biochemical semantics (1/8 requirements)
        \item ODEs: Excellent for metabolism, but no discrete stochastic (4/8 requirements)
        \item Hybrid Petri nets: Close, but lack test/inhibitor arcs and biochemical formulas (4/8 requirements)
        \item Rule-based models: Good for combinatorial complexity, but sequential and discrete-only (2/8 requirements)
    \end{itemize}
    
    \item \textbf{Integration gap}: Current practice uses separate models with manual coupling. Need \textbf{unified formalism}.
\end{enumerate}

\paragraph{The Solution:} Extended Bio-Petri Nets provide the \textbf{first unified formalism} for integrated systems biology, addressing all eight requirements through four core innovations:
\begin{enumerate}
    \item Weak independence \& cooperative parallelism (R5, R8)
    \item Heterogeneous transition types (R1, R2, R6)
    \item Arc-level regulation (R3, R4)
    \item Atomic conservation (R7)
\end{enumerate}

\paragraph{Next Chapters:}
\begin{itemize}
    \item \textbf{Chapter 4}: Extended Bio-PN definition (12-tuple formalism addressing all 8 requirements)
    \item \textbf{Chapter 5}: Weak independence theory (requirement R8, parallelism)
    \item \textbf{Chapter 6}: Biochemical formula tracking (requirement R7, elemental conservation)
\end{itemize}
